\documentclass[10pt,a4paper]{amsart}
\usepackage[margin=19mm]{geometry}
\usepackage{amsmath,amssymb,mathtools}
\usepackage[scr=rsfs,cal=euler]{mathalfa}
\usepackage{xfrac}
\usepackage[unicode,hidelinks,bookmarks=false]{hyperref}
\newcommand{\ie}{i.e.\ }
\newcommand{\cf}{cf.\ }
\newcommand{\resp}{respectively\ }
\newcommand{\Subsection}{\subsection}
\providecommand{\nobreakdash}{\nobreak\hbox{-}}
\setlength{\emergencystretch}{2em}
\multlinegap0pt
\usepackage{bm}
\usepackage{bbm}
\usepackage{dsfont}
\usepackage{xspace}
\usepackage{cancel}
\usepackage{mathtools}
\usepackage{amsthm}
\newtheorem{theorem}{Theorem}
\newtheorem{lemma}[theorem]{Lemma}
\title[Moir\'e Pattern of Rotated Sierpi\'nski Gaskets]{Moir\'e Pattern of Rotated Sierpi\'nski Gaskets}
\author{Ziyu Neroli}
\date{Source: arXiv:2609.30235v1 (CC BY 4.0)}
\begin{document}
\maketitle

\noindent\emph{Source and licence.} Moir\'e Pattern of Rotated Sierpi\'nski Gaskets. Version arXiv:2609.30235v1 (CC BY 4.0), distributed under the Creative Commons Attribution 4.0 International licence, as recorded in the arXiv licence metadata for this exact version and shown on the abstract page https://arxiv.org/abs/2609.30235v1. Full rights notice: licenses/arxiv-2609.30235-CC-BY-4.0.txt.\par\smallskip

\noindent\textit{Editorial scope.} Counted unit: the Lemma whose source label is \texttt{lem:ae-dichotomy} in the article (source0.tex lines 845--848, source label \texttt{lem:ae-dichotomy}), delivered together with the article's own complete original proof of it (source0.tex lines 850--881, one \texttt{proof} environment of 32 lines). No mathematics is written, repaired or re-derived: statement and proof are the article's own text line for line. Required context retained uncounted: lem:ae-exp-sep (lines 806--809). Every definition, display and label of those lines is retained unchanged. Omitted from this package and not redistributed: 1--805, 810--844, 849--849, 882--1209. Nothing inside a retained excerpt is omitted or altered: every retained body is delivered exactly as the article's lines are written, so no omission receipt of any kind accompanies this package. Citations: the retained excerpts carry no citation command at all, so no bibliography entry of the article is transcribed and no reference is redistributed. Reference adaptation: none was needed and none was made; every reference inside the delivered unit resolves inside it, so no unresolved reference or citation is produced. Printed slips kept literal: no printed word, symbol, hypothesis or number of the counted statement or of its proof was altered. Layout and class: the article's own class is \texttt{ctexart} with packages including geometry, amsmath, amssymb, amsthm, mathtools, graphicx, float, xurl, hyperref, tikz, authblk, pgfplots; the wrapper is the lane's pinned \texttt{amsart} editorial wrapper at 10pt on A4, which restates the theorem-like environments the retained text needs and adds the article's own macro definitions that the retained text needs (none), with extra wrapper packages (bm, bbm, dsfont, xspace, cancel, mathtools). The delivered numbering is the unit's own. Assets: no retained span contains a figure environment, a graphics-inclusion command, a PDF-page inclusion, a table or a TikZ picture; no image file is copied into this package, the package declares no asset dependency and no image file is redistributed with it. Licence: version arXiv:2609.30235v1 of this article is distributed under the Creative Commons Attribution 4.0 International licence, as recorded in the arXiv licence metadata for this exact version and shown on the abstract page https://arxiv.org/abs/2609.30235v1. A full rights notice is delivered at licenses/arxiv-2609.30235-CC-BY-4.0.txt. Characters: every retained excerpt is pure ASCII, so the delivered engine, XeLaTeX with its default Latin Modern text font, reports no missing character.
\begin{lemma}
\label{lem:ae-exp-sep}
For Lebesgue-almost every $\theta$, the system $\{\phi_{ij}^{\theta}:i,j\in\mathcal I\}$ satisfies exponential separation.
\end{lemma}

\begin{lemma}
\label{lem:ae-dichotomy}
For Lebesgue-almost every $\theta$, every direction $\vec e\in S^1$ has the following property: at least one of $\mu_{\vec e}$ and $\mu_{\mathbf R_\theta^T\vec e}$ satisfies exponential separation.
\end{lemma}

\begin{proof}
We first exclude angles at which two lattice vectors become too nearly parallel.
For non-zero $\vec\lambda,\vec\eta\in\boldsymbol\Lambda$, the determinant $\vec\lambda\wedge\mathbf R_\theta\vec\eta$ is a sine function of $\theta$ with amplitude $|\vec\lambda||\vec\eta|\ge1$.
Hence the set where $0<|\vec\lambda\wedge\mathbf R_\theta\vec\eta|<\delta$ has length at most $C\delta$.
At level $n$, there are at most $C16^n$ pairs with lattice coordinates bounded by $2^{n+1}$.
The same summation as in Lemma~\ref{lem:ae-exp-sep}, with threshold $16^{-n}n^{-2}$, shows that for almost every $\theta$, every such pair satisfies
\[
\vec\lambda\wedge\mathbf R_\theta\vec\eta=0
\quad\text{or}\quad
|\vec\lambda\wedge\mathbf R_\theta\vec\eta|\ge16^{-n}n^{-2}
\]
at all sufficiently large levels.
Fix such an angle outside $\Theta_0$, so the zero alternative is excluded.

Suppose that both projected systems fail exponential separation for some $\vec e$.
For every $A>0$ and all sufficiently large $n$, their cylinder differences then give non-zero lattice vectors $\vec\lambda_n,\vec\eta_n$, with coordinates bounded by $2^{n+1}$, such that
\[
|\langle\vec e,\vec\lambda_n\rangle|\le e^{-An},
\qquad
|\langle\vec e,\mathbf R_\theta\vec\eta_n\rangle|\le e^{-An}.
\]
Here multiplication of the cylinder differences by $2^n$ is absorbed by choosing a larger exponential rate before rescaling.
Both vectors lie within distance $e^{-An}$ of the line perpendicular to $\vec e$, and their lengths are at most $C2^n$.
Their determinant therefore satisfies
\[
16^{-n}n^{-2}
\le |\vec\lambda_n\wedge\mathbf R_\theta\vec\eta_n|
\le C2^ne^{-An}.
\]
Taking $A=4$ gives $n^{-2}\le C(32e^{-4})^n$, which is impossible for all large $n$.
The exceptional angle set was chosen before $\vec e$, so the conclusion holds simultaneously in every direction.
\end{proof}

\end{document}
